\documentclass[12pt]{article}
\usepackage[letterpaper,top=0.75in,bottom=0.75in,left=0.78in,right=1.55in,marginparwidth=1.15in,marginparsep=0.18in]{geometry}
\usepackage{fontspec}
\setmainfont{Latin Modern Roman}
\usepackage{microtype}
\usepackage{fancyhdr}
\usepackage{titlesec}
\usepackage{enumitem}
\usepackage{xcolor}
\usepackage[hidelinks]{hyperref}
\setlength{\parindent}{1.05em}
\setlength{\parskip}{0.32em}
\setlength{\headheight}{15pt}
\titleformat{\section}{\large\bfseries\scshape}{}{0pt}{}
\titleformat{\subsection}{\normalsize\bfseries}{}{0pt}{}
\pagestyle{fancy}
\fancyhf{}
\fancyhead[L]{Whately Logic: Julius Caesar Lit-Anchor}
\fancyhead[R]{Lesson 4}
\fancyfoot[C]{\thepage}
\renewcommand{\headrulewidth}{0.4pt}

\newcommand{\focusbox}[1]{\par\vspace{0.35em}\noindent\begingroup\setlength{\fboxsep}{6pt}\colorbox{gray!12}{\begin{minipage}{\dimexpr\linewidth-2\fboxsep\relax}#1\end{minipage}}\endgroup\par\vspace{0.45em}}
\newcommand{\studentline}{\par\vspace{0.68em}\hrule\vspace{0.72em}}
\newcommand{\shortline}{\par\vspace{0.42em}\hrule\vspace{0.55em}}
\newcommand{\gloss}[1]{\marginpar{\scriptsize\raggedright\setlength{\baselineskip}{7.8pt}\color{gray!70!black}#1}}
\newcommand{\speaker}[1]{\par\vspace{0.35em}\noindent\textsc{#1}\par\vspace{-0.15em}}

\begin{document}
\begin{center}
{\LARGE\bfseries Julius Caesar Lit-Anchor}\par\vspace{0.18em}
{\large\scshape The Serpent's Egg}\par\vspace{0.35em}
{Lesson 4: Practical Syllogisms, Middle Terms, and Dangerous Conclusions}\par\vspace{0.55em}
\rule{0.82\textwidth}{0.4pt}
\end{center}

\section*{Purpose}
A lit-anchor connects the week's logic lesson to a recurring literary text. In Lesson 4, Whately asks us to see how two claims can be arranged so that they seem to force a conclusion. A syllogism map names the major premise, minor premise, conclusion, major term, minor term, and middle term. The most important question is whether the middle term truly connects the premises to the conclusion.

\textit{Julius Caesar} gives us a serious test case. Brutus does not speak like a fool. He does not merely say, ``I dislike Caesar.'' He turns political fear into an argument. That makes his reasoning more dangerous and more important to examine. A syllogism map can show whether Brutus has proved his conclusion or only made fear look disciplined.

\focusbox{\textbf{Whately's Lesson 4 habit:} Slow an argument down into two premises and a conclusion. Identify the middle term. Then ask whether that middle term really proves the conclusion.}

\section*{Central Question}
Does Brutus's ``serpent's egg'' argument prove that Caesar must be killed, or does it only make fear sound like proof?

\section*{Guided Text: Act 2, Scene 1}
Brutus is alone in his orchard. Caesar has not yet been crowned. Brutus is trying to decide whether Caesar should be killed before he has the chance to become a tyrant. Read the original text below. The side notes are light glosses only; they do not replace the text.

\speaker{Brutus}
\begin{verse}
It must be by his death: and for my part,\\
I know no personal cause to spurn at him,\\
But for the general. He would be crown'd:\\
How that might change his nature, there's the question.\gloss{Brutus admits he has no private grievance. His argument depends on what Caesar might become.}\\
It is the bright day that brings forth the adder;\\
And that craves wary walking. Crown him?--that;--\\
And then, I grant, we put a sting in him,\\
That at his will he may do danger with.\\
The abuse of greatness is, when it disjoins\\
Remorse from power: and, to speak truth of Caesar,\\
I have not known when his affections sway'd\\
More than his reason.\gloss{This is a crucial concession: Brutus has not actually seen Caesar ruled by passion more than reason.} But 'tis a common proof,\\
That lowliness is young ambition's ladder,\\
Whereto the climber-upward turns his face;\\
But when he once attains the upmost round.\\
He then unto the ladder turns his back,\\
Looks in the clouds, scorning the base degrees\\
By which he did ascend. So Caesar may.\\
Then, lest he may, prevent.\gloss{The argument turns on possibility: because Caesar may become dangerous, Brutus says Rome should prevent it.} And, since the quarrel\\
Will bear no colour for the thing he is,\\
Fashion it thus; that what he is, augmented,\\
Would run to these and these extremities:\\
And therefore think him as a serpent's egg\\
Which, hatch'd, would, as his kind, grow mischievous,\\
And kill him in the shell.\gloss{The image makes preemptive killing seem prudent: destroy the danger before it hatches.}
\end{verse}

\section*{Guided Analysis}
Brutus's speech can be emotionally persuasive because it feels cautious and public-spirited. Whately helps us turn the speech into a practical syllogism map.

\subsection*{A possible syllogism map}
\begin{itemize}[leftmargin=*]
\item \textbf{Major premise:} Anyone who will probably become a dangerous tyrant if given supreme power should be stopped before he gains that power.
\item \textbf{Minor premise:} Caesar will probably become a dangerous tyrant if crowned.
\item \textbf{Conclusion:} Therefore, Caesar should be stopped before he is crowned.
\item \textbf{Major term:} people who should be stopped before gaining power.
\item \textbf{Minor term:} Caesar.
\item \textbf{Middle term:} one who will probably become a dangerous tyrant if given supreme power.
\end{itemize}

\subsection*{Where the argument is strongest}
The argument is not nonsense. Brutus is asking a real political question: should a city wait until tyranny is complete before acting against it? His major premise has some force. If a danger is certain and grave, prevention may be wiser than reaction.

\subsection*{Where the argument is weakest}
The weak point is the minor premise. Has Brutus proved that Caesar will probably become a tyrant? Brutus himself says, ``I have not known when his affections sway'd / More than his reason.'' He then shifts from evidence about Caesar to a general claim about ambition: people often become proud when they climb high. That may be true as a warning, but it does not by itself prove Caesar's future guilt.

\focusbox{\textbf{Lesson 4 takeaway:} Brutus's syllogism depends on a middle term: ``future tyrant.'' If Caesar truly belongs under that middle term, the conclusion becomes much stronger. If Caesar has not been proved to belong under it, the argument may turn fear of possibility into permission for violence.}

\section*{Discussion}
Use these questions to prepare for the independent passage.

\begin{enumerate}[leftmargin=*]
\item What conclusion is Brutus trying to justify?
\item What is the difference between proving Caesar \textit{is} dangerous and proving Caesar \textit{may become} dangerous?
\item Which premise is more debatable: the major premise or the minor premise? Why?
\item Why is the ``serpent's egg'' image persuasive?
\item Does the image prove the argument, or does it only make the argument feel more vivid?
\item What would Brutus need to prove before his conclusion became justified?
\end{enumerate}

\section*{Independent Text: Act 1, Scene 2}
Cassius also tries to move Brutus against Caesar. His reasoning is different from Brutus's orchard soliloquy. Cassius argues from Caesar's weakness, Brutus's honor, Roman liberty, and resentment of one-man greatness. Read the passage, then map the argument Cassius implies.

\speaker{Cassius}
\begin{verse}
Why, man, he doth bestride the narrow world\\
Like a Colossus, and we petty men\\
Walk under his huge legs and peep about\\
To find ourselves dishonourable graves.\gloss{Cassius frames Caesar's greatness as Roman humiliation.}\\
Men at some time are masters of their fates:\\
The fault, dear Brutus, is not in our stars,\\
But in ourselves, that we are underlings.\\
Brutus and Caesar: what should be in that ``Caesar''?\\
Why should that name be sounded more than yours?\\
Write them together, yours is as fair a name;\\
Sound them, it doth become the mouth as well;\\
Weigh them, it is as heavy; conjure with 'em,\\
Brutus will start a spirit as soon as Caesar.\gloss{Cassius argues from equality of names and noble reputation toward political equality.}\\
Now, in the names of all the gods at once,\\
Upon what meat doth this our Caesar feed,\\
That he is grown so great? Age, thou art shamed!\\
Rome, thou hast lost the breed of noble bloods!\\
When went there by an age, since the great flood,\\
But it was famed with more than with one man?\\
When could they say till now, that talk'd of Rome,\\
That her wide walls encompass'd but one man?\\
Now is it Rome indeed and room enough,\\
When there is in it but one only man.\\
O, you and I have heard our fathers say,\\
There was a Brutus once that would have brook'd\\
The eternal devil to keep his state in Rome\\
As easily as a king.\gloss{Cassius invokes Brutus's ancestor and Rome's hatred of kings.}
\end{verse}

\section*{Independent Analysis}
Answer in complete thoughts. Leave yourself enough evidence that this work can become a portfolio artifact.

\begin{enumerate}[leftmargin=*]
\item What conclusion does Cassius want Brutus to reach about Caesar?\studentline\studentline
\item Write one possible major premise behind Cassius's argument.\studentline\studentline
\item Write one possible minor premise about Caesar or Brutus.\studentline\studentline
\item What is the middle term that is supposed to connect Cassius's premises to his conclusion?\studentline\studentline
\item Does Cassius prove that Caesar is dangerous, or does he prove only that Caesar is not naturally superior to Brutus? Explain the difference.\studentline\studentline\studentline
\item Which argument is more logically disciplined: Brutus's serpent's egg argument or Cassius's Colossus argument? Why?\studentline\studentline
\end{enumerate}

\section*{Portfolio Artifact: Julius Caesar Practical Syllogism Map}
Choose either Brutus's argument or Cassius's argument. Complete the map below carefully enough that it can be saved in your Logic Portfolio.

\textbf{Original dramatic claim or passage:}\studentline
\textbf{Major premise:}\studentline
\textbf{Minor premise:}\studentline
\textbf{Conclusion:}\studentline
\textbf{Major term:}\studentline
\textbf{Minor term:}\studentline
\textbf{Middle term:}\studentline
\textbf{Does the middle term really connect the premises to the conclusion?}\studentline
\textbf{Which premise is most debatable, and what would the speaker need to prove?}\studentline\studentline

\enlargethispage{3\baselineskip}
\section*{Closing Synthesis}
Whately's practical syllogism map does not tell us automatically whether Brutus is a patriot, a murderer, a tragic idealist, or a political fool. It does something more disciplined: it shows where the argument bears weight. Brutus's conclusion depends on the middle term ``future tyrant.'' Cassius's conclusion depends on a different middle term, something like ``one man falsely treated as greater than Rome's nobles.'' In both cases, Shakespeare shows that political action can be driven by arguments that sound noble before they have been fully proved. Logic asks whether the bridge is strong enough to carry the conclusion.

\end{document}
